\documentclass[11pt]{article}

\usepackage[margin=1in]{geometry}
\usepackage[T1]{fontenc}
\usepackage{lmodern}
\usepackage{amsmath,amssymb}
\usepackage{enumitem}
\usepackage[colorlinks=true,linkcolor=blue,citecolor=blue,urlcolor=blue]{hyperref}

\emergencystretch=3em

\title{Literature Status Note for arXiv:math/0512207}
\author{Run \texttt{fa\_banach\_001}}
\date{2026-07-03}

\begin{document}
\maketitle

\noindent\textbf{Status.}
\texttt{literature\_already\_answered}. This note records a later-literature
answer to the Slicing Problem signal in Emanuel Milman's
\emph{Dual Mixed Volumes and the Slicing Problem}, arXiv:math/0512207. It is
not a new proof from this run.

\section*{Source Question}

The source paper is E. Milman, \emph{Dual Mixed Volumes and the Slicing
Problem}, arXiv:math/0512207.
In the Introduction, pages 1--2 of the arXiv PDF, Milman recalls the
Slicing Problem. In the notation of that paper, the question is whether
there is a universal constant \(C\) such that
\[
  L_K \leq C
\]
for every centrally symmetric convex body \(K \subset \mathbb{R}^n\), after
putting \(K\) in isotropic position and normalizing its volume to one.
Milman also gives Bourgain's equivalent hyperplane-section formulation: every
centrally symmetric volume-one convex body should have an \((n-1)\)-dimensional
section whose volume is bounded below by a universal positive constant.

\section*{Supporting Answer}

The supporting paper is B. Klartag and J. Lehec, \emph{Affirmative Resolution
of Bourgain's Slicing Problem using Guan's Bound}, arXiv:2412.15044.
On page 1, the abstract states that the paper gives the final affirmative
step for Bourgain's slicing problem. The theorem asserted there is that for
every convex body \(K \subseteq \mathbb{R}^n\) of volume one, there exists a
hyperplane \(H\) and a universal constant \(c>0\) such that
\[
  \operatorname{Vol}_{n-1}(K \cap H) > c .
\]
This is precisely the hyperplane-section formulation recalled in the source
paper, and it is stronger in scope than the centrally symmetric case mentioned
there. Therefore the source's Slicing Problem signal is now answered by later
literature.

\section*{Why This Is Not A New Result}

The supporting authors explicitly identify their theorem as an affirmative
resolution of Bourgain's slicing problem. The relation is therefore an exact
later-literature answer, not an agent-derived proof, counterexample, or
hidden implication.

\section*{Scope Limitations}

This packet records only the literature status of the Slicing
Problem/hyperplane conjecture signal in arXiv:math/0512207. It does not address
the separate question in the source about whether \(\mathrm{BP}_k^n =
\mathrm{I}_k^n\). It also does not address the later radial-sum questions or
the outer mean-radius assumptions for the cube and for unconditional bodies.
Those remain separate mathematical targets.

\section*{Search And Verification Notes}

The run indexes were searched for the source arXiv id, title terms, and core
keywords including \texttt{slicing}, \texttt{hyperplane conjecture},
\texttt{dual mixed volumes}, \texttt{BP\_k}, and \texttt{radial sums}. No exact
packet for arXiv:math/0512207 was found. A related prior packet for a different
source, arXiv:1310.1204, already recorded arXiv:2412.15044 as the decisive
answer to Bourgain's slicing problem; its supporting PDF is copied into this
packet.

\section*{Local Files}

\begin{itemize}[leftmargin=2em]
  \item Source PDF: \texttt{source\_paper.pdf}.
  \item Supporting PDF: \texttt{supporting\_paper\_2412.15044.pdf}.
  \item Ledger record:
    \texttt{0512207\_slicing\_answered\_by\_2412.15044.json}.
\end{itemize}

\begin{thebibliography}{9}
\bibitem{Milman2005}
E. Milman,
\emph{Dual Mixed Volumes and the Slicing Problem},
arXiv:math/0512207.

\bibitem{KlartagLehec2024}
B. Klartag and J. Lehec,
\emph{Affirmative Resolution of Bourgain's Slicing Problem using Guan's Bound},
arXiv:2412.15044.
\end{thebibliography}

\end{document}
